% TODO make hw stylesheet
\documentclass[10pt]{article}
\usepackage[letterpaper,top=1in]{geometry}
\usepackage[dvipsnames,svgnames]{xcolor}
\usepackage[shortlabels]{enumitem}
\usepackage[framemethod=TikZ]{mdframed}
\usepackage{amsmath,amssymb,amsthm}
\usepackage[colorlinks]{hyperref}
\usepackage{multicol}
\usepackage{tabularx}

\hypersetup{
	colorlinks=true,
	linkcolor=blue,
	filecolor=magenta,
	urlcolor=magenta,
	pdftitle={MAT 528 HW}
}

\usepackage{mathtools}
\usepackage{thmtools}
\usepackage[textsize=footnotesize]{todonotes}
\usepackage{titling}

\reversemarginpar
\mdfdefinestyle{mdbluebox}{roundcorner=10pt,innerbottommargin=9pt,
	linecolor=blue,backgroundcolor=TealBlue!5,}
\declaretheoremstyle[headfont=\sffamily\bfseries\color{MidnightBlue},
mdframed={style=mdbluebox},]{thmbluebox}

\mdfdefinestyle{mdbloobox}{frametitlefont=\bfseries,innerbottommargin=8pt,
	nobreak=true,backgroundcolor=TealBlue!5,linecolor=blue,}
\declaretheoremstyle[headfont=\bfseries\color{MidnightBlue},
mdframed={style=mdbloobox},headpunct={\\[3pt]},postheadspace=0pt,]{thmbloobox}

\mdfdefinestyle{mdredbox}{frametitlefont=\bfseries,innerbottommargin=8pt,
	nobreak=true,backgroundcolor=Salmon!5,linecolor=RawSienna,}
\declaretheoremstyle[headfont=\bfseries\color{RawSienna},
mdframed={style=mdredbox},headpunct={\\[3pt]},postheadspace=0pt,]{thmredbox}

\mdfdefinestyle{mdgreenbox}{linecolor=ForestGreen,backgroundcolor=ForestGreen!5,
	linewidth=2pt,rightline=false,leftline=true,topline=false,bottomline=false,}
\declaretheoremstyle[headfont=\bfseries\sffamily\color{ForestGreen!70!black},
mdframed={style=mdgreenbox},headpunct={ --- },]{thmgreenbox}

\mdfdefinestyle{mdorangebox}{linecolor=RedOrange,backgroundcolor=RedOrange!5,
	linewidth=2pt,rightline=false,leftline=true,topline=false,bottomline=false,}
\declaretheoremstyle[headfont=\bfseries\sffamily\color{RedOrange!70!black},
mdframed={style=mdorangebox},headpunct={ --- },]{thmorangebox}

\mdfdefinestyle{mdblackbox}{linecolor=black,backgroundcolor=RedViolet!5!gray!5,
	linewidth=3pt,rightline=false,leftline=true,topline=false,bottomline=false,}
\declaretheoremstyle[mdframed={style=mdblackbox}]{thmblackbox}

\declaretheoremstyle[spaceabove=6pt,spacebelow=6pt]{basehead}

\declaretheorem[style=basehead,name=Problem]{problem}
\declaretheorem[style=thmredbox,name=Problem,numbered=no]{problem*}
\declaretheorem[style=thmbloobox,name=Setup]{setup} % for config geo
\declaretheorem[style=thmredbox,name=Example,sibling=problem]{example}
\declaretheorem[style=thmredbox,name=$\bigstar$ Problem,sibling=problem]{niceprob}
\declaretheorem[style=thmbluebox,name=Theorem,sibling=problem]{theorem}
\declaretheorem[style=thmbluebox,name=Corollary,sibling=theorem]{corollary}
\declaretheorem[style=thmbluebox,name=Proposition,sibling=theorem]{prop}
\declaretheorem[style=thmbluebox,name=Lemma,numberwithin=problem]{lemma}
\declaretheorem[style=thmbluebox,name=Theorem,numbered=no]{theorem*}
\declaretheorem[style=thmbluebox,name=Lemma,numbered=no]{lemma*}
\declaretheorem[style=thmgreenbox,name=Claim,numberwithin=problem]{claim}
\declaretheorem[style=thmgreenbox,name=Claim,numbered=no]{claim*}
\declaretheorem[style=thmblackbox,name=Remark,numberwithin=problem]{remark}
\declaretheorem[style=thmblackbox,name=Remark,numbered=no]{remark*}
\declaretheorem[style=thmgreenbox,name=Definition,sibling=theorem]{definition}
\declaretheorem[style=thmgreenbox,name=Definition,numbered=no]{definition*}

\providecommand{\ol}{\overline}
\providecommand{\ceil}[1]{\left\lceil #1 \right\rceil}
\providecommand{\floor}[1]{\left\lfloor #1 \right\rfloor}
\newcommand{\dd}{\mathrm{d}}
\providecommand{\ul}{\underline}
\providecommand{\eps}{\varepsilon}
\providecommand{\half}{\frac{1}{2}}
\providecommand{\CC}{\mathbb C}
\providecommand{\FF}{\mathbb F}
\providecommand{\NN}{\mathbb N}
\providecommand{\QQ}{\mathbb Q}
\providecommand{\RR}{\mathbb R}
\providecommand{\ZZ}{\mathbb Z}
\providecommand{\dg}{^\circ}
\providecommand{\ii}{\item}
\providecommand{\setdiff}{\smallsetminus}
\providecommand{\alert}[1]{{\sffamily\textbf{\textcolor{blue}{#1}}}}
\providecommand{\inv}{^{-1}}
\providecommand{\mb}[1]{\mathbf{#1}}

\newenvironment{soln}{\begin{proof}[Solution]}{\end{proof}}

\providecommand{\pow}{\operatorname{Pow}}
\providecommand{\vocab}[1]{{\textbf{\textcolor{ForestGreen}{#1}}}}
\providecommand{\lcm}{\operatorname{lcm}}
\providecommand{\abs}[1]{\left|#1\right|}

\usepackage{mathrsfs}

% place title at top right
\setlength{\droptitle}{-8ex}
\pretitle{\begin{flushright}\Large\bfseries}
\posttitle{\par\end{flushright}}
\preauthor{\begin{flushright}}
\postauthor{\end{flushright}}
\predate{\begin{flushright}}
\postdate{\end{flushright}}

\title{MAT 528 HW02}
\author{\large Aditya Pahuja}
\date{\large Due 09/11}
\begin{document}
\maketitle

\begin{problem}
    [Munkres 17.6b]
    Let $A$ and $B$ denote subsets of a space $X$.
    Prove that $\ol{A\cup B} = \ol A\cup \ol B$.
\end{problem}

\begin{soln}
    First, $\ol{A\cup B}$ is a closed set containing $A$ and $B$,
    so it contains both $\ol A$ and $\ol B$,
    hence $\ol A\cup\ol B\subseteq\ol{A\cup B}$.
    Second, $\ol{A}\cup\ol B$ is a closed set containing
    both $A$ and $B$, so it contains $A\cup B$ and thus contains its closure,
    i.e. $\ol{A\cup B}\subseteq\ol A\cup\ol B$.

    We have shown that $\ol{A\cup B}$ and $\ol A\cup\ol B$ contain each other,
    so the two sets are equal.
\end{soln}

\begin{problem}
    [Munkres 17.11]
    Show that the product of two Hausdorff spaces is Hausdorff.
\end{problem}

\begin{soln}
    Let $X$ and $Y$ be Hausdorff spaces and let
    $(x_1,y_1)$ and $(x_2,y_2)$ be two distinct points
    in the product space $X\times Y$.
    Then, at least one of $x_1\ne x_2$ and $y_1\ne y_2$ holds.

    If the former holds, then since $X$ is Hausdorff
    we can pick disjoint open sets
    $U_1\ni x_1$ and $U_2\ni x_2$ in $X$;
    then $U_1\times Y$ and $U_2\times Y$ are two
    disjoint open sets in $X\times Y$ (because they are products of open sets)
    that respectively contain $(x_1,y_1)$ and $(x_2,y_2)$.

    On the other hand, if the latter holds, then since $Y$ is Hausdorff
    we can pick disjoint open sets
    $V_1\ni y_1$ and $V_2\ni y_2$ in $Y$;
    then $X\times V_1$ and $X\times V_2$ are two
    disjoint open sets in $X\times Y$ (because they are products of open sets)
    that respectively contain $(x_1,y_1)$ and $(x_2,y_2)$.

    Therefore, regardless of which two points $(x_1,y_1)$ and $(x_2,y_2)$
    in $X\times Y$ are chosen, we can find disjoint open sets in $X\times Y$
    separating the two points.
\end{soln}

\begin{problem}
    [Munkres 17.19a]
    If $A\subseteq X$, we define the \vocab{boundary} of $A$ by the equation
    \begin{equation*}
	\partial A = \ol A\cap\ol{(X\setminus A)}.
    \end{equation*}
    Show that $\operatorname{Int}(A)$ and $\partial A$ are disjoint,
    and $\ol A = \operatorname{Int}(A)\cup\partial A$.
\end{problem}

\begin{soln}
    We claim that $\operatorname{Int}(A)$ is the complement of
    $\ol{(X\setminus A)}$. 

    If $x\in\operatorname{Int}(A)$, then there is an open subset
    $U\subseteq X$ such that $x\in U\subseteq A$.
    This means that $U$ does not intersect $X\setminus A$,
    so $x$ cannot be in $\ol{(X\setminus A)}$, i.e.
    $\operatorname{Int}(A)$ is a subset of the complement of
    $\ol{(X\setminus A)}$.

    Conversely, if $x$ is in the complement of $\ol{(X\setminus A)}$,
    then $x\in A$ and there exists an open set $U$ containing $x$
    which does not intersect $X\setminus A$, so $U\subseteq A$;
    the existence of such a $U$ is what it means for $x$ to be in
    $\operatorname{Int}(A)$, hence the complement of $\ol{(X\setminus A)}$
    is a subset of the interior of $A$.

    We have shown inclusion in both directions, so the claim is proven.

    Now, $\operatorname{Int}(A)$ and $\partial A$ must be disjoint
    because the latter is a subset of the complement
    $X\setminus \operatorname{Int}(A) = \ol{(X\setminus A)}$,
    and furthermore
    \begin{equation*}
	\ol A = (\ol A \cap\ol{(X\setminus A)})
	\cup (\ol A \cap(X\setminus \ol{(X\setminus A)}))
	= \partial A \cup (\ol A \cap \operatorname{Int}(A))
	= \partial A \cup \operatorname{Int}(A)
    \end{equation*}
    where the last equality is a result of
    $\operatorname{Int}(A)\subseteq A\subseteq \ol A$.
\end{soln}

\pagebreak

\begin{problem}
    [Munkres 17.20e]
    Find the boundary and interior of the set
    \begin{equation*}
	E = \{(x,y) : 0 < x^2-y^2 \le 1\}\subseteq\RR^2.
    \end{equation*}
\end{problem}

\begin{soln}
    Define $f\colon\RR^2\to\RR$ by $f(x,y) = x^2 - y^2$,
    which is continuous.
    Then, I claim that $F = f\inv([0,1])$ (which is a closed set containing $E$
    by continuity) is the closure of $E$.
    The only thing to check is that the points in $f\inv(0)$
    are all limit points of $E$, since the rest of the points
    in $F$ are in $E$ already. If $x^2-y^2=0$
    and $U$ is an open set containing $(x,y)$, then $U$ contains
    some open ball $B_\eps((x,y))$ centered at $(x,y)$ with $0 < \eps < 1$.
    Let $r = \min(\eps/(1000x), \eps/1000)$. Then, if $x$ is positive,
    \begin{equation*}
	0 < (x+r)^2 - y^2 = (x+r)^2 - x^2 = 2rx + r^2
	\le \frac{\eps}{500} + \frac{\eps^2}{1000^2} < \eps < 1,
    \end{equation*}
    and if $x$ is negative we instead bound $(x-r)^2 - y^2$ in the same way,
    so this ball centered at $(x,y)$ intersects $E$ at a point besides $(x,y)$
    (either $(x-r,y)$ or $(x+r,y)$) making $(x,y)$ a limit point of $E$.

    By similar reasoning we can conclude that the points in $f\inv(1)$
    are also limit points of $\RR^2\setminus E$, meaning any open ball
    around one of these points will contain a point outside $E$.
    Hence $f\inv((0,1))$ is the largest open set contained in $E$,
    so it is the interior of $E$; thus the boundary of $E$ is
    $f\inv([0,1])\setminus f\inv((0,1)) = f\inv(\{0,1\})$.
\end{soln}

\begin{problem}
    [Munkres 18.2]
    Suppose $f\colon X\to Y$ is continuous. If $x$ is a limit point of
    the subset $A$ of $X$, is it necessarily true that
    $f(x)$ is a limit point of $f(A)$?
\end{problem}

\begin{soln}
    No, this statement may not hold; for example,
    consider $f\colon\RR\to\RR$ given by $f(x) = 17$ for all $x$.
    This is continuous because the pre-image of an open set in $\RR$
    is either empty (if the set doesn't contain $17$)
    or all of $\RR$ (if the set does contain $17$).
    Every point of $\RR$ is a limit point of $\RR$ essentially trivially,
    but no point of $\RR$ is a limit point of $f(\RR) = \{1\}$:
    since $\{1\}$ is closed, the only possible limit point is $1$,
    and no ball centered at $1$ intersects $\{1\}$ at a point other than $1$.
\end{soln}

\begin{problem}
    [Munkres 18.5]
    Show that the subspace $(a,b)$ of $\RR$ is homeomorphic with $(0,1)$
    and the subspace $[a,b]$ of $\RR$ is homeomorphic with $[0,1]$.
\end{problem}

\begin{soln}
    Take $f\colon(a,b)\to(0,1)$ defined by $f(x) = \frac{x-a}{b-a}$.
    This is clearly continuous, and we can compute that its inverse is
    $f\inv(y) = (b-a)y + a$, which is also continuous,
    so $f$ is a homeomorphism between $(a,b)$ and $(0,1)$.

    We can define $g\colon[a,b]\to[0,1]$ in the same way
    as $g(x) = \frac{x-a}{b-a}$; the inverse is again
    the continuous map $g\inv(g) = (b-a)y + a$, so $g$ is also a homeomorphism
    between $[a,b]$ and $[0,1]$.
\end{soln}

\begin{problem}
    [Munkres 18.12]
    Let $F\colon\RR^2\to\RR$ be defined by the equation
    \begin{equation*}
        F(x,y) = \begin{cases}
	    xy/(x^2+y^2) & (x,y)\ne(0,0)\\
	    0 & (x,y) = (0,0).
        \end{cases}
    \end{equation*}
    \begin{enumerate}[(a)]
        \ii Show that $F$ is continuous in each variable separately.
	\ii Compute the function $g\colon\RR\to\RR$
	defined by $g(x) = F(x,x)$.
	\ii Show that $F$ is not continuous.
    \end{enumerate}
\end{problem}

\begin{soln}
    (a) The function is symmetric in $x$ and $y$,
    so we only need to prove this for $x$. If we fix $y=y_0\ne 0$, then
    $f_{y_0}(x) = F(x,y_0) = \frac{xy_0}{x^2+y_0^2}$ is obviously continuous
    on $\RR$ because its denominator never vanishes.
    If $y_0=0$ on the other hand, then $f_{y_0}(x)=0$ for all $x$,
    which is trivially continuous.

    (b) $g(x) = F(x,x) = \frac{x^2}{x^2+x^2} = \half$.

    (c) The set $f\inv(\{\half\})$ is the pre-image of a closed set,
    so it should be closed. However, $f(x,y) = \half$ if and only if
    $(x,y)\ne0$ and $x^2+y^2 = 2xy\iff (x-y)^2 = 0\iff x=y$, which
    means that the pre-image is the set $\{(x,x) : x\in\RR,\; x\ne 0\}$,
    which isn't closed because it has $(0,0)$ as a limit point
    (since $B_\eps((0,0))$ contains $(\eps/2,\eps/2)$ for all $\eps>0$)
    but it doesn't contain $0$.
\end{soln}










\end{document}

